\section{The TCDA Framework}\label{sec:framework}

Building on the preliminaries of \Cref{sec:preliminaries}, we now propose
\emph{Topological Causal Data Analysis} (TCDA) as a coherent research
domain at the intersection of TDA and causal inference. Our exposition
proceeds in five steps. \Cref{sec:formal-tcda} gives a formal definition
of a TCDA problem as a tuple linking a topological space of observations
to a causal graph through a functorial topological summary, together with
two foundational propositions establishing well-posedness.
\Cref{sec:domain-A,sec:domain-B,sec:domain-C,sec:domain-D} then defend a
four-domain taxonomy: representation (A), robustness (B), interpretation
(C), and methodology (D), each with three sub-domains. The taxonomy is
summarised in \Cref{fig:taxonomy} and is the paper's primary intellectual
contribution. Every sub-domain is anchored to at least one published
work, ensuring that the framework is grounded rather than aspirational.

\subsection{Formal definition of a TCDA problem}\label{sec:formal-tcda}

The unifying intuition is that, in a TCDA problem, causal inference is
performed not on raw observations but on \emph{topological summaries} of
those observations, where the summary is computed by a functor that
satisfies the stability and interpretability properties recalled in
\Cref{sec:prelim_tda}. The causal question of interest is then either
(i)~a query about the underlying causal structure that generated the
observations, or (ii)~a contrast between topological summaries induced
by different interventions, or (iii)~a robustness/identifiability
statement about classical causal estimands whose validity is controlled
by topological invariants.

\begin{definition}[TCDA problem]\label{def:tcda}
A \emph{topological causal data analysis problem} is a tuple
\[
\TCDA \;=\; (\Mc,\, G,\, \Tc,\, \Cc),
\]
whose components are:
\begin{enumerate}[leftmargin=2em,label=(\roman*)]
\item $\Mc$ is a (Polish) topological space of observations equipped with
a metric or filtration parameter $\rho:\Mc\times\Mc\to\R_{\geq 0}$, on which
probability measures and their interventional counterparts are well-defined;
\item $G = (V,E)$ is a causal structure---typically a DAG or ADMG---whose
vertices represent (possibly latent) random variables $X = (X_v)_{v\in V}$
taking values in $\Mc$ and whose edges encode direct causal relationships
in the SCM sense of \Cref{sec:prelim_ci};
\item $\Tc \,:\, \mathbf{Filt}(\Mc) \,\to\, \mathbf{Pers}$ is a
\emph{topological functor} from a category of filtered objects on $\Mc$
to a category of persistence-like invariants. Canonical instances are
$\Tc = \PHk{k}$ for $k\in\{0,1,2\}$, $\Tc = \Reeb_f$, $\Tc = \Map_f$,
and their functional summaries $\Tc = \varphi(\cdot;\,r)$ (silhouettes)
or $\Tc = \Lambda$ (landscapes). We require $\Tc$ to be stable with
respect to some interleaving-type metric on $\mathbf{Filt}(\Mc)$;
\item $\Cc$ is a \emph{causal query}---a functional of the
interventional or counterfactual distribution induced by $G$ on $\Mc$.
We allow $\Cc$ to act directly on raw outcomes ($\Cc = \ATE,\CATE,\TATE,
P(y\mid \doop(x)),\ldots$) or on topological summaries
($\Cc = W_p(\Dgmk{k}(P^{x}),\Dgmk{k}(P^{x'}))$, for example, as in
\citet{kim2026topological}).
\end{enumerate}
A TCDA \emph{solution} is an estimator $\widehat\Cc_n$ together with
calibrated uncertainty (confidence sets, hypothesis tests, posterior
credible sets) for the population quantity $\Cc$, computed from a sample
$\{X_i\}_{i=1}^n \subset \Mc$ using $\Tc$ as a (possibly sufficient)
statistic.
\end{definition}

\Cref{def:tcda} is intentionally broad. The four components
$(\Mc, G, \Tc, \Cc)$ play orthogonal roles: $\Mc$ specifies \emph{what
the data is}, $G$ specifies \emph{which causal relationships are
admitted}, $\Tc$ specifies \emph{how shape is summarised}, and $\Cc$
specifies \emph{the scientific question}. Existing work occupies
distinct corners of this space:
\begin{itemize}[leftmargin=2em]
\item \citet{ibeling2021topological} take $\Tc$ to be the identity and
endow $\Mc = \mathfrak{M}$ (a space of SCMs) with a topology that controls
the Pearl--Bareinboim causal hierarchy.
\item \citet{kim2026topological} fix $\Cc$ to be a Wasserstein contrast
between persistence diagrams of potential outcomes and develop AIPW
estimators with $\sqrt n$-rates and weak-convergence guarantees.
\item \citet{elyaagoubi2023statistical} fix $\Mc$ to be a space of
multivariate time-series, $\Tc = \PHk{1}$ of a windowed dependence
network, and $\Cc$ to be a hypothesis test about the number of cycles.
\item \citet{lecci2014statistical} develops the calibrated statistical
inference for $\Tc$ alone that any TCDA estimator must inherit.
\end{itemize}
Thus the four seed works of the field each fix three components of the
TCDA tuple and study the remaining one. The TCDA programme is precisely
the systematic study of the joint behaviour of all four.

% =====================================================================
% Figure: TCDA Taxonomy (Section 3.1)
% This is a working draft. The user has indicated that final figure
% production will be handled separately; the file is provided so the
% document compiles and the cross-reference \Cref{fig:taxonomy} resolves.
% =====================================================================
\begin{figure}[t]
\centering
\begin{tikzpicture}[
  font=\small,
  level distance=11mm,
  level 1/.style={sibling distance=42mm, level distance=12mm,
                  every node/.style={draw, rounded corners, fill=blue!8,
                                     align=center, minimum width=32mm,
                                     minimum height=7mm, font=\small\bfseries}},
  level 2/.style={sibling distance=14mm, level distance=13mm,
                  every node/.style={draw, rounded corners, fill=teal!6,
                                     align=center, minimum width=28mm,
                                     minimum height=6mm, font=\footnotesize}},
  edge from parent/.style={draw, -{Latex[length=2mm]}, semithick}]

\node[draw, rounded corners, fill=violet!10, minimum width=64mm,
      minimum height=9mm, font=\bfseries] {Topological Causal Data Analysis (TCDA)}
  child {node {A.\,Topological\\Causal Representation}
    child {node {A1\\PH for Causal\\Structure Learning}}
    child {node {A2\\Mapper \& Reeb\\for Abstraction}}
    child {node {A3\\Causal Embeddings\\in Persistence}}
  }
  child {node {B.\,Topological\\Robustness}
    child {node {B1\\Stability of\\Causal Estimates}}
    child {node {B2\\Topological\\Confounding}}
    child {node {B3\\Functorial\\Causal Inference}}
  }
  child {node {C.\,Causal Interpretation\\of Topological Features}
    child {node {C1\\Betti Numbers as\\Causal Explanation}}
    child {node {C2\\Persistent Cohomology\\for Mediation}}
    child {node {C3\\Stratified Causal\\Heterogeneity}}
  }
  child {node {D.\,Applications\\\& Methodology}
    child {node {D1\\Time-Series\\Causality}}
    child {node {D2\\Spatial \&\\Spatiotemporal}}
    child {node {D3\\Industry\\Applications}}
  };
\end{tikzpicture}
\caption{The proposed TCDA taxonomy. Four orthogonal domains, each with
three sub-domains, organise the field. Domain A concerns
\emph{representation}: how topology summarises the structure of causal
data. Domain B concerns \emph{robustness}: how topological stability
controls the sensitivity of causal estimates to perturbations. Domain C
concerns \emph{interpretation}: which causal stories are encoded by
topological features. Domain D concerns \emph{methodology and applications}
in concrete data modalities. The taxonomy is intended to be exhaustive
(every existing TCDA contribution we are aware of fits in at least one
node) but not mutually exclusive (many works contribute to several
sub-domains simultaneously).}
\label{fig:taxonomy}
\end{figure}


\paragraph{Requirements on $\Tc$.}
For \Cref{def:tcda} to be useful in a causal context, the functor $\Tc$
must satisfy three operating conditions:
\begin{enumerate}[leftmargin=2em,label=(F\arabic*)]
\item \emph{Stability.} Small perturbations of the input filtration
should produce small perturbations of $\Tc(\cdot)$ in an interleaving-type
metric. Persistent homology satisfies this with constant $1$ for the
bottleneck distance (\Cref{eq:stability}). Mapper and Reeb functors
admit weaker but useful stability properties under appropriate
discretisation \citep{munch2016convergence,carriere2018statistical}.
\item \emph{Interpretability.} The features of $\Tc(\cdot)$ should admit
a substantive interpretation in the underlying scientific domain (e.g.,
$1$-cycles in $\PHk{1}$ of a brain functional network correspond to
re-entrant cortical loops; \citealp{elyaagoubi2023statistical,
sizemore2018cliques}).
\item \emph{Computability.} The functor must be computable in time
polynomial in the sample size for the kind of data the application
demands. Off-the-shelf libraries \texttt{Ripser}, \texttt{GUDHI},
\texttt{giotto-tda} provide $O(n^3)$ persistent-homology computation up
to dimension $2$ in practice; sparse and approximate variants extend
this to $n \gtrsim 10^5$ \citep{bauer2021ripser,otter2017roadmap,
tauzin2021giotto}.
\end{enumerate}
We refer to functors satisfying (F1)--(F3) as \emph{admissible} and
restrict attention to admissible $\Tc$ in all that follows.

We now state two propositions that establish the well-posedness of the
TCDA tuple. The first transfers the celebrated stability of persistent
homology to any Lipschitz causal query; the second pins down the
informational condition under which a causal graph is identifiable from
a topological functor alone. Both are stated for the canonical case
$\Tc = \PHk{k}$ but extend immediately to silhouettes, landscapes, and
other Lipschitz functional summaries.

\begin{proposition}[Stability transfer]\label{prop:stability-transfer}
Let $\TCDA = (\Mc, G, \Tc, \Cc)$ be a TCDA problem with $\Tc = \PHk{k}$
applied to the Vietoris--Rips filtration of a compact subset of $\Mc$,
and let
\[
  \Cc\bigl[\Tc(P), \Tc(Q)\bigr] \;=\; \phi\bigl(\Dgmk{k}(P),\,\Dgmk{k}(Q)\bigr)
\]
for some functional $\phi$ that is $L$-Lipschitz with respect to the
$p$-Wasserstein distance on persistence diagrams. Then for any pair of
distributions $P, P' \in \mathcal{P}(\Mc)$ with bounded support and any
fixed $Q$,
\[
\bigl|\Cc[\Tc(P),\Tc(Q)] - \Cc[\Tc(P'),\Tc(Q)]\bigr|
\;\leq\; L \cdot W_p\!\bigl(\Dgmk{k}(P),\,\Dgmk{k}(P')\bigr)
\;\leq\; L \cdot d_H\!\bigl(\supp P,\, \supp P'\bigr),
\]
where $d_H$ denotes the Hausdorff distance between the supports of $P$
and $P'$.
\end{proposition}

\begin{proof}[Sketch]
The first inequality is the Lipschitz assumption on $\phi$; the second
is the stability theorem for persistent homology
(\Cref{eq:stability}; \citealp{cohen2007stability,chazal2016structure}).
The constant is $1$ in the canonical bottleneck case and tracks the
$p$-Wasserstein constant otherwise.
\end{proof}

\emph{Significance.} \Cref{prop:stability-transfer} says that any causal
query expressible as a Lipschitz functional of persistence diagrams
inherits the quantitative robustness of $\PHk{k}$: a small Hausdorff
perturbation of the support---a much weaker requirement than convergence
in any parametric metric---bounds the perturbation of the causal
estimand. This delivers a robustness guarantee with no analogue in
classical (parametric) causal inference. \citet{kim2026topological}
exploit precisely this transfer to establish stability of their TATE
estimator under the silhouette-based summary $\Tc = \varphi(\cdot;r)$.

\begin{proposition}[Identifiability from a topological functor]
\label{prop:identifiability}
Let $G$ vary over a class $\mathcal{G}$ of DAGs over the random vector
$X$, let $\mathcal{P}_G$ denote the observational distribution induced by
$G$, and let $\Tc$ be an admissible topological functor. Then $G$ is
identifiable from $\Tc(\mathcal{P}_G)$ alone---i.e., there is a function
$\Psi$ such that $\Psi(\Tc(\mathcal{P}_G)) = G$ for all $G \in \mathcal{G}$
\emph{if and only if} the assignment $G \mapsto \Tc(\mathcal{P}_G)$ is
injective on $\mathcal{G}$.
\end{proposition}

\begin{proof}[Sketch]
The ``only if'' direction is immediate: if the map is non-injective,
then two distinct $G_1 \neq G_2$ produce identical $\Tc(\cdot)$, so no
inverse $\Psi$ exists. The ``if'' direction is the definition of
identifiability: an injective map admits a left inverse on its image,
and any such left inverse plays the role of $\Psi$. Standard measurable
selection (e.g., the Kuratowski--Ryll-Nardzewski theorem) supplies a
measurable $\Psi$ when $\mathcal{G}$ is countable, as in the
finite-DAG case.
\end{proof}

\emph{Significance.} \Cref{prop:identifiability} isolates the exact
informational condition under which a topology-based discovery algorithm
can hope to recover the underlying graph. It allows existing
topology-based discovery procedures to be classified by what $\Tc$ they
implicitly use: TOPIC \citep{tran2024constraining} uses the
score-Jacobian functor, SSTS \citep{xu2025information} uses a mutual-
information-based functor, and Hodge-decomposition methods
\citep{nocurl2024} use the harmonic component of a $1$-chain. The
proposition does not, of course, certify that the injectivity condition
\emph{holds} for a given $\Tc$; we list this as
\Cref{q:tcda-identifiability} in \Cref{sec:research}.

\medskip
Together, \Cref{prop:stability-transfer,prop:identifiability} establish
that the TCDA tuple is non-vacuous: it inherits both the stability and
the identifiability properties that motivate its definition. The
remaining sections of this section show that the four domains of the
taxonomy fall out naturally as the four ``faces'' of \Cref{def:tcda}---one
for each of the components $\Tc$, $\Cc$ in their representational and
robustness roles, plus one for the interpretive direction (reading
causality \emph{out} of $\Tc$) and one for the methodological direction
(applying TCDA in specific data modalities).

\subsection{Domain A: Topological Causal Representation}\label{sec:domain-A}

Domain A addresses the question: \emph{which topological structure best
represents the causal content of high-dimensional, non-Euclidean data?}
The role of $\Tc$ in \Cref{def:tcda} is here that of a sufficient
statistic for the causal structure or for the interventional contrasts
of interest. Three sub-domains arise.

\subsubsection*{A1. Persistent homology for causal structure learning.}
Persistent homology of a multivariate distribution or its empirical
dependence graph can serve as a sufficient statistic for the underlying
causal DAG. Two threads converge here. First, modern \emph{topological-
ordering} approaches to causal discovery---e.g.\ TOPIC
\citep{tran2024constraining}, SSTS \citep{xu2025information}, NoCurl-X
\citep{nocurl2024}, OptiTopo \citep{wu2026optimization}, and CaPS
\citep{rolland2022score}---use rank-based or information-theoretic
functors to extract a topological order from the data and then refine
edges within that order; the order itself is a topological invariant of
the joint distribution. Second, persistent homology of the
\emph{dependence network} of a multivariate time series identifies
feedback cycles \citep{elyaagoubi2023statistical}; loops in
$\PHk{1}$ correspond to directed feedback loops in the underlying
Granger graph, and a generative simulation framework
\citep{elyaagoubi2023statistical} now exists to test how well
TDA-based discovery methods recover known cycle structure. In the
language of \Cref{prop:identifiability}, sub-domain A1 chooses $\Tc$ to
satisfy the injectivity condition for restricted classes of $G$ (e.g.\
additive-noise models with non-Gaussian innovations).

\subsubsection*{A2. Mapper and Reeb graphs for causal abstraction.}
Where A1 produces a graph that approximates $G$, A2 produces a
\emph{topological abstraction} of the underlying causal mechanism: a
simplification of the data space along a filter function chosen to
align with the causal structure. The Mapper algorithm
\citep{singh2007mapper,carriere2018statistical} partitions $\Mc$ along
level sets of a filter $f$ and glues the resulting clusters into a
simplicial complex; under appropriate choices of $f$, the resulting
complex captures coarse causal regimes (treatment-response strata,
phase transitions, mechanism-shift points). Reeb graphs play the
analogous role for continuous filters \citep{biasotti2008reeb,
morozov2013interleaving}. The recent ball-mapper variant
\citep{dlotko2019ball} is particularly attractive for causal abstraction
because it bypasses the cover-selection problem and admits a stability
guarantee. The statistical scaffolding of \citet{lecci2014statistical}
applies directly: bootstrap confidence sets on cluster trees yield
calibrated abstractions, which is essential when the abstraction is to
serve as the carrier of a downstream causal query.

\subsubsection*{A3. Causal embeddings in persistence diagrams.}
Sub-domain A3 represents the \emph{causal effect itself} as a
topological feature: each potential outcome distribution $P^a$ is
summarised by a persistence diagram $\Dgmk{k}(P^a)$ (or its silhouette
$\varphi(\cdot;P^a,r)$) and the causal contrast is a distance,
hypothesis test, or learned mapping between these diagrams. This is
exactly the formulation of \citet{kim2026topological}, who define the
\emph{topological average treatment effect}
\[
\TATE = W_1\!\bigl(\Dgmk{k}(P^{a=1}),\,\Dgmk{k}(P^{a=0})\bigr),
\]
and provide an AIPW estimator with $\sqrt n$-rate. The same idea applies
in non-Euclidean outcome spaces: the geodesic ATE of
\citet{kurisu2024geodesic} for Riemannian-valued outcomes admits a
topological generalisation in which the Fr\'echet mean is replaced by
the persistence diagram. Sub-domain A3 thus exhibits the TCDA tuple in
its most direct form: $\Tc$ converts outcomes to diagrams; $\Cc$ is a
metric or test on those diagrams; the causal graph $G$ underwrites
identifiability of the contrast.

\subsection{Domain B: Topological Robustness in Causal Inference}\label{sec:domain-B}

Domain B addresses the question: \emph{when can topological invariants
\textbf{control} the sensitivity of causal estimates to perturbations of
the data, the model, or the identifying assumptions?} Here the role of
$\Tc$ is to deliver Lipschitz-type bounds that propagate to $\Cc$.

\subsubsection*{B1. Stability of causal estimates under topological perturbations.}
\Cref{prop:stability-transfer} is the prototype: any Lipschitz causal
query whose primary statistic is a persistence diagram inherits a
quantitative robustness bound from the bottleneck/Wasserstein stability
of $\PHk{k}$. \citet{kim2026topological} exploit this in their stability
theorem for the silhouette TATE; \citet{farzam2025topology} extend it to
representation-learning settings where topological summaries serve as
distributionally-robust losses. A complementary thread uses persistence
landscapes \citep{bubenik2015statistical} as smooth, Hilbert-space-valued
embeddings, which allows classical M-estimation rates to be transferred
to TDA-based causal estimators. The statistical inference programme of
\citet{lecci2014statistical}---confidence sets, bootstrap, rates of
convergence on landscapes and diagrams---is the inferential backbone of
all B1 work: a TCDA estimator that lacks calibrated uncertainty cannot
deliver causal robustness, only descriptive plots.

\subsubsection*{B2. Topological detection of unobserved confounding.}
A persistent ``hole'' in the data manifold can betray the action of an
unobserved variable: when an unobserved confounder $U$ moves an outcome
along an orbit, the support of the observed distribution exhibits a
non-trivial $\PHk{1}$ class that the observed graph cannot explain. This
observation has been operationalised in two ways. First, \emph{anomaly
scores} based on bottleneck or Wasserstein distance flag samples whose
local persistence diagram differs from the population diagram
\citep{bruel2018topological,nigmetov2022topological}; flagged samples
are candidates for hidden-confounder action. Second, the topological
hierarchy of \citet{ibeling2021topological} shows that the set of SCMs
for which an assumption-free identification of $P(y\mid \doop(x))$ is
possible is \emph{topologically meager}, formalising the intuition that
confounders are ``everywhere dense.'' This puts a precise upper limit on
what B2 detection methods can hope to achieve: they identify residual
topology consistent with confounding but cannot, in general, certify the
absence of confounding without further assumptions.

\subsubsection*{B3. Functorial causal inference.}
Sub-domain B3 asks for the categorical analogue of \Cref{def:tcda}: a
formal description of how causal inference and topological filtration
interact as composable functors. The commutative square of
\Cref{fig:functorial} captures the ideal case: applying $\Tc$ to a pair
of interventional distributions and then taking a topological contrast
should yield the same answer as taking the causal contrast at the level
of distributions and then applying $\Tc$. \citet{mahadevan2025sheaf,
mahadevan2021causal,mahadevan2022universal} have developed a sheaf-
theoretic framework in which causal inference is presented as a
universal construction in a topos of sheaves over a base category of
contexts; do-calculus is recovered through the subobject classifier.
This programme can be read as a B3 formalisation of TCDA when $\Tc$ is
chosen to be the global-section functor of an appropriate sheaf. The
Alexandroff-topology approaches of \citet{willig2025alexandroff}
provide a more elementary entry point.

% =====================================================================
% Figure: Functorial commutative square for TCDA (Section 3.3 / 3.5)
% Working draft -- the user will refine final visual styling.
% =====================================================================
\begin{figure}[t]
\centering
\begin{tikzcd}[column sep=large, row sep=large]
P^{X = x}
  \arrow[r, "\Tc"] \arrow[d, "\doop(x') / \doop(x)"', shift right=2pt]
  \arrow[d, "T_{\!\Cc}"', shift left=10pt, dashed]
  & \Tc(P^{X=x})
  \arrow[d, "\Phi_{\Cc}", shift left=2pt]
  \\
P^{X = x'}
  \arrow[r, "\Tc"']
  & \Tc(P^{X=x'})
\end{tikzcd}
\caption{Functorial picture of a TCDA query. The horizontal arrows apply
a topological functor $\Tc$ (e.g.\ $\PHk{k}$, $\Reeb$, $\Map$) to two
interventional distributions. The left vertical arrow $\doop(x)\!\to\!\doop(x')$
acts at the level of the causal model; the right vertical arrow $\Phi_{\Cc}$
acts at the level of topological summaries. A TCDA solution requires either
that the diagram \emph{commutes} -- so that the topological summary of a
causal contrast equals the causal contrast of topological summaries -- or
that the failure to commute is itself a quantifiable, interpretable causal
object. Stability of $\Tc$ (cf.\ \Cref{eq:stability}) makes the right
vertical arrow Lipschitz with respect to bottleneck distance; hence under
commuting the left vertical arrow inherits a Lipschitz bound in the same
metric.}
\label{fig:functorial}
\end{figure}


\subsection{Domain C: Causal Interpretation of Topological Features}\label{sec:domain-C}

Domain C inverts the direction of arrows: rather than using topology to
build or stabilise causal estimates, it reads \emph{causal stories} off
topological features that have already been computed. The role of $\Tc$
remains as before, but the role of $\Cc$ is now that of an interpretive
overlay---a translation of a topological summary into the language of
mechanisms, mediators, and heterogeneity.

\subsubsection*{C1. Betti numbers as causal explanations.}
A persistent $k$-cycle is an invariant; whether it admits a causal
interpretation depends on the data-generating process. In multivariate
time-series neuroscience, a robust $1$-cycle in the Granger-network
filtration of a healthy brain corresponds to a re-entrant cortical loop
that the disease modifies \citep{sizemore2018cliques,
fathian2022trend,elyaagoubi2023statistical}. In gene-regulation,
$1$-cycles in correlation-network filtrations have been linked to
feedback motifs whose perturbation alters phenotype
\citep{kovacev2016using,cang2018integration}. In each case, the
topological feature is a candidate causal mechanism whose validity is
adjudicated by intervention. Sub-domain C1 develops the principled
language for these moves: when is a persistent cycle ``causal'' rather
than ``spurious''?

\subsubsection*{C2. Persistent cohomology for mediation analysis.}
Cohomology is more naturally suited to decomposition than is homology:
a coboundary represents the action of a potential function on a chain,
and the Hodge decomposition of a flow on a simplicial complex separates
the gradient (direct), curl (feedback), and harmonic (global)
components \citep{nocurl2024,desilva2007coverage}. We argue that this
decomposition is the right topological tool for \emph{mediation
analysis}: the total causal effect of an exposure $A$ on an outcome $Y$
through mediators $M$ decomposes naturally into a gradient component
(direct effect), a curl component (cyclic feedback through $M$), and a
harmonic component (effects not attributable to any local mediator).
\Cref{q:cohomology-mediation} in \Cref{sec:research} states the
corresponding open problem.

\subsubsection*{C3. Stratified spaces and causal heterogeneity.}
Real causal effects are often heterogeneous across strata defined not
by explicit covariates but by latent topological features: disease
subtypes are connected components of a phenotype manifold; phase
transitions are codimension-$1$ singularities. The theory of stratified
spaces \citep{hughes2004stratified,banagl2007topological} provides a
framework for causal effects that vary across topological strata,
generalising the conditional ATE to a stratum-conditional ATE indexed
by a topological feature. The cluster-tree statistical inference of
\citet{lecci2014statistical} and the topological clustering theory of
\citet{chazal2013persistence} provide the inferential machinery; the
causal content remains to be developed and is one of the main open
avenues we identify in \Cref{sec:research}.

\subsection{Domain D: Methodology and Applications}\label{sec:domain-D}

Domain D is the applied face of TCDA. It does not introduce new
mathematical primitives---these are supplied by Domains A--C---but it
identifies the data modalities and substantive areas in which the
TCDA tuple is most readily instantiated. Each application sub-domain
specialises the tuple $(\Mc, G, \Tc, \Cc)$ to a class of problems.

\subsubsection*{D1. TCDA for time-series causality.}
Time-series data are the natural habitat for sub-domain D1. The space
$\Mc$ is a phase-space reconstruction (Takens embedding) of one or more
observed series; $G$ encodes temporal directed causality (Granger or
its generalisations); $\Tc = \PHk{1}$ of the reconstructed phase-space
detects periodicity and feedback structure; $\Cc$ is a hypothesis test
or estimation of dynamical causality. \citet{bando2022causal} combine
persistent homology of Takens embeddings with convergent cross mapping;
\citet{harnack2017causal,gholizadeh2018short} use persistence to
augment Granger-type tests. \citet{elyaagoubi2023statistical} provide a
simulation framework on which D1 methods can be benchmarked, with
prescribed numbers of cycles serving as ground truth.

\subsubsection*{D2. TCDA for spatial and spatiotemporal causality.}
The space $\Mc$ is a manifold ($\R^2$, $\R^3$, a sphere) or a
spatio-temporal product; $G$ encodes spatially-distributed causal
mechanisms (e.g.\ teleconnections in climate, propagation in epidemics);
$\Tc$ is typically cubical persistent homology
\citep{kaczynski2004computational} on a gridded scalar field. The
cosmological literature \citep{sousbie2011persistent,
vandeweygaert2011alpha,cisewski2014nonparametric} provides a mature
template for descriptive TCDA on cosmic-web data; sub-domain D2 is the
research programme that adds causal attribution to such descriptions
(e.g.\ attribution of an extreme-precipitation event to a Pacific
oscillation pattern via TCDA contrast). \citet{kim2026topological}'s
machinery extends naturally here when the outcome of interest is a
gridded field rather than a vector.

\subsubsection*{D3. Industry applications: healthcare, finance, climate.}
Sub-domain D3 collects the application-driven contributions that fall
outside the cleaner D1/D2 modalities. Healthcare TCDA includes
topological biomarkers of disease progression and treatment response
\citep{kovacev2016using,cang2018integration,iqbal2025classification,
soares2020sars,kim2018topological}, where \citet{kim2026topological}
exhibit the SARS-CoV-2 CT-scan application and the GEOM-Drugs molecular
graph application. Finance TCDA includes topological early-warning
indicators for systemic risk \citep{gidea2018topological,
ismail2022early}, where the causal content is the attribution of a
regime change to a triggering market factor. Climate TCDA includes
topological characterisations of extreme-event clusters
\citep{cisewski2014nonparametric,strommen2025topological}, where the
causal content is the attribution of regime change to forcing. In each
case the TCDA tuple specialises predictably; the open methodological
problem is the calibration of $\widehat\Cc_n$ under realistic
non-stationarity.

\subsection{Reading the framework}\label{sec:reading}

\Cref{fig:taxonomy} is intended to be read in two ways. \emph{Vertically},
each domain A--D represents a different way to fix the TCDA tuple: A
fixes $\Tc$ and $\Cc$ to representational roles; B fixes them to
robustness roles; C inverts the direction by asking what
$\Cc$ \emph{means} given $\Tc$; D specialises $\Mc$ and $G$ to
particular data modalities. \emph{Horizontally}, the sub-domains within
each row form a methodological progression: from the simplest tool
(homology/diagram) at left to the most structural (functorial/
stratified) at right.

The framework is exhaustive in the following operational sense: every
work we cite in this paper occupies at least one node of
\Cref{fig:taxonomy}, and most occupy more than one---reinforcing the
claim that the field is real but currently fragmented. The framework is
not, however, mutually exclusive; nor is it final. \Cref{sec:literature}
populates the taxonomy with the actual state of the literature, and
\Cref{sec:research} returns to it to formulate the open questions and
research avenues that each node generates.

We close this section with an observation that motivates the rest of
the paper. Each of the four domains was identified by a single seed
work: \citet{ibeling2021topological} essentially establishes B2--B3 (the
topology of SCMs and the meagre set of identifiable models);
\citet{kim2026topological} essentially establishes A3--B1 (causal
contrasts as topological functionals with stability); the statistical
programme of \citet{lecci2014statistical} provides the inferential
machinery that A--C all require; and the time-series framework of
\citet{elyaagoubi2023statistical} anchors D1 and supplies the
simulation infrastructure on which the rest of the methodology can be
benchmarked. The TCDA programme thus begins with four seed-points and
the conjecture, defended in the remainder of this paper, that they are
the corners of a single coherent research domain.
